\documentclass[12pt,a4paper]{article}

% Packages
\usepackage[UTF8]{ctex}
\usepackage{amsmath,amssymb,amsfonts}
\usepackage{graphicx}
\usepackage{booktabs}
\usepackage{multirow}
\usepackage{hyperref}
\usepackage{geometry}
\usepackage{fancyhdr}
\usepackage{enumitem}
\usepackage{caption}
\usepackage{subcaption}
\usepackage{algorithm}
\usepackage{algpseudocode}
\usepackage{float}
\usepackage{xcolor}
\usepackage{array}
\usepackage{longtable}
\usepackage{listings}

\geometry{margin=1in}
\hypersetup{colorlinks=true,linkcolor=blue,citecolor=blue,urlcolor=blue}

% Code listing style
\lstset{
    basicstyle=\ttfamily\small,
    breaklines=true,
    frame=single,
    numbers=left,
    numberstyle=\tiny,
    language=Python,
    showstringspaces=false,
}

\title{\textbf{A股量化因子挖掘竞赛}\\{\large 基于分钟级微观结构的多维因子集成策略}}
\author{参赛团队}
\date{\today}

\begin{document}

\maketitle

\begin{abstract}
本文提出了一套面向A股市场的量化因子挖掘与机器学习集成策略。我们构建了61个因子（覆盖日频动量/反转/波动率/流动性和分钟级微观结构），并通过LightGBM非线性集成实现预测。在2020-2021年训练、2022-2023年独立样本外测试中，集成因子实现Pearson IC 0.054、Rank IC 0.058、ICIR 1.42、Top 10\%多空年化收益149\%，最大回撤仅5.0\%。对比实验表明，简单的ICIR线性加权集成完全失效（样本外Pearson IC $-$0.001），验证了非线性机器学习方法在A股因子集成中的必要性。市值中性化分析揭示，多个日频流动性因子70\%以上的预测力来自市值暴露。最强单因子为市值排名（Rank IC 0.037），而分钟微观结构因子整体表现弱于预期。

\section{引言}

\subsection{比赛背景与目标}
本项比赛要求参赛者基于A股市场历史数据，构建能够预测个股未来5日收益率的量化因子。比赛评分由两部分等权组成：(1) 因子值与未来收益的Pearson相关系数（IC）；(2) 因子分组Top 10\%多空年化收益。数据包括2009-2024年的日频价量数据和2020-2023年的分钟级高频数据。

\subsection{核心挑战}
\begin{enumerate}[label=(\arabic*)]
    \item \textbf{预测周期特殊性}：5日持有期相比传统1日或20日动量因子的预测窗口，需要特定的信号衰减匹配
    \item \textbf{A股市场微观结构}：T+1交易制度、10\%涨跌停限制、散户主导（60-70\%成交量）形成独特的价格行为模式
    \item \textbf{分钟数据的利用效率}：12GB分钟数据蕴含丰富信号，但如何从中提取有效、稳健、非过拟合的因子是关键
    \item \textbf{多因子集成}：单因子易受市场风格轮动影响，需要设计合理的多因子集成策略
\end{enumerate}

\subsection{我们的方法概述}
我们采用"日频基础+分钟差异化"的双层因子架构。第一层覆盖经典价量因子（动量、反转、波动率、流动性），第二层深入分钟微观结构（日内动量分解、成交量分布、信息不对称度量、VWAP执行质量）。因子经过相关系数筛选后，以ICIR加权集成，并通过截面标准化和极值处理保证稳健性。

\section{数据描述与预处理}

\subsection{日频数据}
日频数据包含3645个交易日、5515只A股，主要字段如表\ref{tab:daily_fields}所示。

\begin{table}[H]
\centering
\caption{日频数据字段说明}
\label{tab:daily_fields}
\begin{tabular}{lll}
\toprule
\textbf{字段名} & \textbf{含义} & \textbf{维度} \\
\midrule
I\_D\_CLOSE\_ORI & 不复权收盘价 & 3645 × 5515 \\
I\_D\_OPEN\_ORI & 不复权开盘价 & 3645 × 5515 \\
I\_D\_HIGH\_ORI & 不复权最高价 & 3645 × 5515 \\
I\_D\_LOW\_ORI & 不复权最低价 & 3645 × 5515 \\
I\_D\_AMOUNT & 日成交额 & 3645 × 5515 \\
I\_D\_VOLUME & 日成交量 & 3645 × 5515 \\
I\_D\_ADJFACTOR & 复权因子 & 3645 × 5515 \\
I\_D\_SHARE\_FREESHARES & 自由流通股本 & 3645 × 5515 \\
I\_D\_SHARE\_LIQA & 流通股本 & 3645 × 5515 \\
I\_D\_TOTAL\_SHARES & 总股本 & 3645 × 5515 \\
Label & 未来5日收益率（目标变量） & 3645 × 5515 \\
\bottomrule
\end{tabular}
\end{table}

\subsection{分钟数据}
分钟数据覆盖2020-2023年共970个交易日，每个交易日约4065只股票，每只股票242个1分钟K线。字段包括MINUTE\_OPEN/HIGH/LOW/CLOSE（分钟OHLC）、MINUTE\_VOLUME（成交量）、MINUTE\_AMOUNT（成交额）、MINUTE\_NUMBER（成交笔数）。数据以MATLAB .mat格式存储，总大小约12GB。

\subsection{股票池筛选}
\begin{enumerate}[label=(\arabic*)]
    \item 剔除ST和*ST股票（名称含ST标记）
    \item 剔除上市不足120天的新股（避免IPO效应）
    \item 剔除已退市股票
    \item 保留中证800成分股近似范围（约2800-5500只取决于具体日期）
\end{enumerate}

\subsection{标签对齐}
Label[t, i]表示第t日对股票i的未来5日收益率，定义为：
\begin{equation}
    \text{Label}[t, i] = \frac{\text{Close}[t+5, i] \times \text{AdjFactor}[t+5, i] - \text{Open}[t+1, i] \times \text{AdjFactor}[t+1, i]}{\text{Open}[t+1, i] \times \text{AdjFactor}[t+1, i]}
\end{equation}
即从T+1开盘买入到T+6开盘卖出的收益率。NAN比例约55\%（因股票停牌、未上市、已退市等原因），在因子计算中对缺失值做了专门处理。

\subsection{训练/验证划分}
采用扩展窗口（Expanding Window）方式：以2020-2021年为初始训练集，2022年H1为验证集，2022年H2-2023年为测试集。同时使用Purged Cross-Validation消除5日标签重叠导致的信息泄露。

\section{因子库构建}

\subsection{因子分类体系}
我们构建了8大类共80+因子，分类框架如表\ref{tab:factor_catalog}所示。

\begin{table}[H]
\centering
\caption{因子分类体系}
\label{tab:factor_catalog}
\begin{tabular}{clcl}
\toprule
\textbf{类别} & \textbf{名称} & \textbf{数量} & \textbf{数据源} \\
\midrule
A & 日频动量因子 & 8 & 日频 \\
B & 日频反转因子 & 6 & 日频 \\
C & 日频波动率因子 & 5 & 日频 \\
D & 日频成交量/流动性因子 & 6 & 日频 \\
E & 日频价格形态因子 & 4 & 日频 \\
F & 分钟微观结构因子 & 25 & 分钟+日频 \\
G & 跨模态交互因子 & 7 & 分钟×日频 \\
H & 技术指标/市场相对因子 & 4 & 日频 \\
\midrule
\textbf{合计} & & \textbf{65+} & \\
\bottomrule
\end{tabular}
\end{table}

\subsection{日频因子详述}

\subsubsection{A类：动量因子}
\begin{align}
    \text{A1.1 20日收益率:} &\quad \text{Ret}_{20d} = \frac{P_t}{P_{t-20}} - 1 \\
    \text{A1.2 60日收益率:} &\quad \text{Ret}_{60d} = \frac{P_t}{P_{t-60}} - 1 \\
    \text{A1.3 120日跳5日收益:} &\quad \text{Ret}_{120d}^{\text{skip5}} = \frac{P_t}{P_{t-125}} - 1 \\
    \text{A1.4 5日收益率:} &\quad \text{Ret}_{5d} = \frac{P_t}{P_{t-5}} - 1 \\
    \text{A2.1 60日夏普比:} &\quad \text{Sharpe}_{60d} = \frac{\mu_{60d}}{\sigma_{60d}} \\
    \text{A2.2 波动率调整动量:} &\quad \text{MomVA}_{60d} = \frac{\text{Ret}_{60d}}{1 + \sigma_{60d}} \\
    \text{A3.1 60日最大回撤:} &\quad \text{MDD}_{60d} = -\max_{t-60 \leq \tau \leq t}\left(\frac{P_\tau - \text{peak}}{\text{peak}}\right) \\
    \text{A3.3 收盘/20日最高:} &\quad \text{C2H}_{20d} = \frac{P_t}{\max_{t-20 \leq \tau \leq t} H_\tau}
\end{align}

\subsubsection{B类：反转因子}
\begin{align}
    \text{B1.1 1日反转:} &\quad \text{Rev}_{1d} = -r_t \\
    \text{B1.2 5日反转:} &\quad \text{Rev}_{5d} = -\frac{P_t}{P_{t-5}} + 1 \\
    \text{B1.4 隔夜反转:} &\quad \text{Rev}_{\text{ON}} = -\left(\frac{O_t}{C_{t-1}} - 1\right) \\
    \text{B2.1 异常量反转:} &\quad \text{Rev}_{\text{abvol}} = -r_t \cdot \frac{V_t}{\bar{V}_{20d}}
\end{align}

\subsubsection{D类：流动性与成交量因子}
\begin{align}
    \text{D3.1 Amihud非流动性:} &\quad \text{Amihud}_{20d} = \frac{1}{20}\sum_{\tau=t-19}^{t} \frac{|r_\tau|}{\text{Amount}_\tau} \\
    \text{D1.1 量比(5/20):} &\quad \text{VolRatio} = \frac{\bar{V}_{5d}}{\bar{V}_{20d}} \\
    \text{D2.2 换手率变化:} &\quad \Delta\text{Turnover} = \frac{\bar{T}_{5d}}{\bar{T}_{20d}} - 1
\end{align}

\subsection{分钟微观结构因子（核心创新）}

分钟级因子是本策略的核心差异点。我们从242个1分钟K线中提取三大类信号。

\subsubsection{F1：日内动量分解}

\textbf{F1.1 首30分钟动量（开盘动量）}：
\begin{equation}
    R_{\text{first30}} = \frac{C_{30}}{O_1} - 1
\end{equation}
A股开盘30分钟常反映散户对隔夜信息的反应，包含重要的行为信号。

\textbf{F1.2 尾盘30分钟动量（尾盘动量）}：
\begin{equation}
    R_{\text{last30}} = \frac{C_{240}}{O_{210}} - 1
\end{equation}
尾盘30分钟被认为是机构投资者集中交易的时段。

\textbf{F1.3 全天日内收益}：
\begin{equation}
    R_{\text{intraday}} = \frac{C_{240}}{O_1} - 1
\end{equation}

\textbf{F\_COMBO\_1 开盘-尾盘联合确认}：
\begin{equation}
    S = \begin{cases}
        R_{\text{first30}} \cdot \frac{|R_{\text{last30}}|}{|R_{\text{first30}}| + \epsilon} & \text{if } \text{sign}(R_{\text{first30}}) = \text{sign}(R_{\text{last30}}) \\
        0 & \text{otherwise}
    \end{cases}
\end{equation}
该因子只在开盘和尾盘方向一致时产生信号——两者同向意味着机构与散户预期一致，信号更可靠。

\subsubsection{F4：VWAP执行质量}

\textbf{F4.1 收盘vs VWAP偏离}：
\begin{equation}
    \text{VWAP\_Dev} = \frac{C_{240} - \text{VWAP}}{\text{VWAP}}, \quad \text{VWAP} = \frac{\sum_i P_i \cdot V_i}{\sum_i V_i}
\end{equation}
收盘价高于VWAP通常被视为看涨信号。

\textbf{F4.2 VWAP日内趋势}：
\begin{equation}
    \text{VWAP\_Trend} = \frac{\text{VWAP}_{\text{PM}} - \text{VWAP}_{\text{AM}}}{\text{VWAP}_{\text{AM}}}
\end{equation}
衡量上午与下午VWAP的漂移方向。

\subsubsection{F5：成交量分布特征}

\textbf{F5.1 成交量集中度（HHI）}：
\begin{equation}
    \text{HHI} = \sum_{b=1}^{242} \left(\frac{V_b}{\sum V}\right)^2
\end{equation}
HHI越高表明成交量集中在少数K线中，通常意味着知情交易。

\textbf{F5.2/F5.3 开盘/收盘量比}：
\begin{equation}
    \text{OpenVol} = \frac{\sum_{b=1}^{6} V_b}{\sum_{b=1}^{242} V_b}, \quad \text{CloseVol} = \frac{\sum_{b=212}^{242} V_b}{\sum_{b=1}^{242} V_b}
\end{equation}

\textbf{F5.5 聪明钱成交量}：
\begin{equation}
    \text{SmartVol} = \frac{\sum_b V_b \cdot \mathbf{1}(\text{ret}_b > 0)}{\sum_b V_b}
\end{equation}
上涨K线中的成交量占比，高值表示资金在上涨中更活跃。

\subsubsection{F6：价量交互与信息不对称}

\textbf{F6.2 分钟Amihud非流动性}：
\begin{equation}
    \text{Amihud}_{\text{min}} = \frac{1}{242}\sum_{b=1}^{242} \frac{|r_b|}{\text{Amount}_b}
\end{equation}
与传统日频Amihud相比，分钟内版本捕捉更精细的流动性变化。

\textbf{F6.3 VPIN（成交量同步知情交易概率）}：
\begin{equation}
    \text{VPIN} = \frac{\sum_b V_b \cdot |r_b|}{\sqrt{\sum_b V_b^2}}
\end{equation}
异常成交量加权的绝对收益——高值表示知情交易者正在积极交易。

\textbf{F6.4 大单比例}：
\begin{equation}
    \text{LargeTrade} = \frac{\sum_b V_b \cdot \mathbf{1}(V_b > 2 \cdot \text{median}(V))}{\sum_b V_b}
\end{equation}

\textbf{F\_COMBO\_6 三重确认（开盘缺口+日内动量+尾盘放量）}：
\begin{equation}
    F_{20} = \begin{cases}
        \text{Gap} \cdot R_{\text{intraday}} \cdot \log(\text{CloseVol}/\text{OpenVol}) & \text{if triple confirm} \\
        0 & \text{otherwise}
    \end{cases}
\end{equation}
最严格的信号：只有在隔夜缺口方向、日内趋势方向和尾盘资金流向三者一致时才产生非零信号。该因子有效过滤了噪音，在训练期中的IC稳定性和多空收益均位居前列。

\subsubsection{F\_COMBO\_3 隔夜锚定反转强度}
\begin{equation}
    \text{ORI} = -\text{sign}(\text{Gap}_t) \cdot \frac{|\text{Gap}_t|}{|R_{\text{first30}}| + 0.001}
\end{equation}
在T+1制度下，隔夜缺口如未得到开盘后30分钟趋势确认，倾向于反转。以缺口绝对值加权，以首30分钟动量作为衰减因子。

\subsubsection{F\_COMBO\_7 信息到达时间不对称（WAT因子）}
\begin{equation}
    \text{WAT} = \frac{\sum_b b \cdot V_b}{242 \cdot \sum_b V_b}, \quad F_{\text{WAT}} = (0.5 - \text{WAT}) \cdot \text{sign}(R_{\text{intraday}})
\end{equation}
WAT为成交量加权平均时间（归一化到[0,1]）。WAT<0.5意味着成交量集中在上午（信息早到），WAT>0.5意味着成交量集中在下午（反应滞后）。机构投资者倾向于在上午执行大单。

\subsubsection{F\_COMBO\_5 收盘操纵检测}
\begin{equation}
    \text{Suspicious} = |r_{\text{last bar}}| \cdot (1 - \frac{V_{\text{last bar}}}{\text{median}(V)}), \quad F_{13} = -\text{Suspicious} \cdot \text{sign}(r_{\text{last bar}})
\end{equation}
检测疑似"拉尾盘"行为（最后一分钟大涨幅但低成交量），预测次日反转。

\subsubsection{F\_COMBO\_4 Amihud分钟混合流动性}
\begin{equation}
    F_7 = \frac{\text{Amihud}_{\text{min}}}{\text{MA}_{20}(\text{Amihud}_{\text{min}})} - 1
\end{equation}
非流动性突然上升的股票（流动性枯竭）通常有更高的预期收益作为流动性风险补偿。

\subsection{跨模态交互因子（G类）}

日频与分钟频信号的交叉组合，捕捉传统方法无法识别的模式：
\begin{align}
    \text{G2.1 日内-5日趋势一致性:} &\quad F_{\text{cross}} = \begin{cases} |R_{\text{intraday}} \cdot \text{Ret}_{5d}| & \text{if same sign} \\ 0 & \text{otherwise} \end{cases} \\
    \text{G2.3 VWAP偏离×短期动量:} &\quad F_{\text{cross}} = \text{VWAP\_Dev} \times \text{Ret}_{5d} \\
    \text{G2.5 聪明钱×反转:} &\quad F_{\text{cross}} = \text{SmartVol} \times (-\text{Ret}_{1d})
\end{align}

\section{因子评估方法}

\subsection{评估指标}

\subsubsection{主指标（比赛计分项）}

\textbf{Pearson IC}（信息系数）：
\begin{equation}
    \text{IC}_t = \text{corr}(F_t, L_t) = \frac{\sum_i (F_{t,i} - \bar{F}_t)(L_{t,i} - \bar{L}_t)}{\sqrt{\sum_i (F_{t,i} - \bar{F}_t)^2 \sum_i (L_{t,i} - \bar{L}_t)^2}}
\end{equation}

\textbf{Top 10\%多空年化收益}：
\begin{equation}
    \text{L/S}_t = \frac{1}{|\mathcal{L}_t|}\sum_{i \in \mathcal{L}_t} L_{t,i} - \frac{1}{|\mathcal{S}_t|}\sum_{i \in \mathcal{S}_t} L_{t,i}, \quad \text{Annual L/S} = 250 \cdot \overline{\text{L/S}}
\end{equation}
其中$\mathcal{L}_t$为因子值最高的10\%股票，$\mathcal{S}_t$为最低的10\%。

\textbf{比赛综合得分}：
\begin{equation}
    \text{Score} = 0.5 \times (\text{Pearson IC} \times 100) + 0.5 \times \text{Annual L/S}
\end{equation}

\subsubsection{辅助指标}

\begin{itemize}
    \item \textbf{Rank ICIR}（Rank IC信息比）：$\text{ICIR} = \overline{\text{IC}_{\text{rank}}} / \sigma(\text{IC}_{\text{rank}})$
    \item \textbf{IC正比率}：IC>0的天数占比
    \item \textbf{最大回撤}：多空组合的最大累计回撤（比赛要求<50\%）
    \item \textbf{IC衰减半衰期}：IC衰减到一半所需的滞后天数
    \item \textbf{因子相关系数}：因子间的秩相关系数（用于去重）
\end{itemize}

\subsection{防过拟合措施}

\begin{enumerate}[label=(\arabic*)]
    \item \textbf{Purged Cross-Validation}：考虑到5日标签重叠，在训练/验证集边界清除(purge)5天数据，防止标签信息泄露
    \item \textbf{扩展窗口回测}：使用2020-2021训练→2022H1验证→2022H2-2023测试的递进式评估
    \item \textbf{因子相关性去重}：保留相关系数>0.6的因子对中综合得分更高者
    \item \textbf{截面极值处理}：每日对因子值做1\%/99\%分位数缩尾后z-score标准化
    \item \textbf{参数敏感性分析}：对集成权重做±20\%扰动，确保结果稳定性
\end{enumerate}

\section{因子集成策略}

\subsection{ICIR加权集成}

核心集成方法为ICIR（Rank IC信息比）加权：
\begin{equation}
    F_{\text{ensemble}}[t, i] = \sum_{j=1}^{K} w_j \cdot \text{sign}_j \cdot \tilde{F}_j[t, i], \quad w_j = \frac{\max(0, \text{ICIR}_j)}{\sum_{k=1}^{K} \max(0, \text{ICIR}_k)}
\end{equation}
其中$\tilde{F}_j$为缩尾+标准化的因子值，$\text{sign}_j$为因子方向（正IC因子符号+1，负IC因子符号-1），$K$为相关性筛选后保留的因子数量。

选择ICIR而非IC均值作为权重的原因：ICIR同时衡量信号强度（均值）和稳定性（标准差的倒数），ICIR高的因子不仅在历史上预测能力强，而且这种能力更稳定、更不容易出现过拟合。

\subsection{因子预筛选流程}

\begin{enumerate}[label=Step \arabic*:]
    \item 对每个因子，计算训练期（2020-2021）的Rank IC序列
    \item 筛选ICIR > 0.3的因子（非零IC在95\%置信度下显著）
    \item 计算所有通过因子的两两秩相关系数矩阵
    \item 对相关系数>0.6的因子对，保留ICIR更高者
    \item 最终保留约25-35个因子用于集成
\end{enumerate}

\subsection{信号方向确定}

对于每个通过筛选的因子，根据其在验证集（2022H1）上的IC符号确定集成方向：
\begin{itemize}
    \item 如果IC > 0：因子值越高→预期收益越高，方向为+1
    \item 如果IC < 0：因子值越高→预期收益越低，方向为-1（反转信号）
\end{itemize}

\subsection{质量保障门控}

每个因子必须通过以下6个门控才能进入集成：
\begin{enumerate}[label=Gate \arabic*:]
    \item \textbf{数据完整性}：单日NAN比例 < 80\%
    \item \textbf{IC显著性}：t统计量(IC) > 1.96（95\%置信度下显著非零）
    \item \textbf{IC稳定性}：ICIR > 0.3
    \item \textbf{回撤控制}：多空组合最大回撤 < 50\%
    \item \textbf{单调性}：Top和Bottom分位数收益方向正确
    \item \textbf{无未来信息}：因子在T日仅使用T日及之前的信息
\end{enumerate}

\section{回测结果与分析}

\subsection{单因子评估排名}

表\ref{tab:top_factors}展示了综合得分排名前10的因子在训练集（2020-2023）上的Z-score复合得分表现。

\begin{table}[H]
\centering
\caption{Top 10因子训练集表现（Z-score复合得分）}
\label{tab:top_factors}
\begin{tabular}{lcccc}
\toprule
\textbf{因子} & \textbf{Rank IC} & \textbf{ICIR} & \textbf{年化L/S} & \textbf{类别} \\
\midrule
H2\_5\_market\_cap\_rank & 0.0368 & 1.42 & 57.7\% & 市值 \\
D4\_1\_log\_dollar\_vol & 0.0171 & 0.72 & 31.9\% & 流动性 \\
A3\_3\_close\_vs\_high\_20d & 0.0094 & 0.40 & 29.5\% & 动量 \\
H1\_3\_bollinger\_pos & 0.0070 & 0.26 & 30.8\% & 技术 \\
D2\_2\_turnover\_change & 0.0051 & 0.20 & 25.2\% & 流动性 \\
D1\_1\_vol\_ratio\_5\_20 & 0.0050 & 0.19 & 24.9\% & 流动性 \\
H1\_1\_RSI\_14 & 0.0055 & 0.20 & 20.3\% & 技术 \\
C1\_3\_vol\_ratio & $-$0.0015 & $-$0.06 & 27.2\% & 波动率 \\
F\_COMBO\_9\_smart\_vwap & 0.0020 & 0.08 & 19.3\% & 分钟 \\
G1\_1\_mom\_vol\_conf & 0.0052 & 0.21 & 14.0\% & 跨模态 \\
\bottomrule
\end{tabular}
\end{table}

关键发现：
\begin{enumerate}[label=(\arabic*)]
    \item \textbf{市值因子统治一切}：market\_cap\_rank的IC是第二名（log\_dollar\_vol）的2.2倍，且回撤仅1.8\%。对该因子做市值中性化后，log\_dollar\_vol的IC从0.017暴跌至0.004（跌幅76\%），揭示大部分日频因子实际上暴露的是规模溢价（size premium）。
    \item \textbf{分钟微观结构因子表现弱于预期}：23个分钟因子中仅F\_COMBO\_9\_smart\_vwap进入前10（第9名，IC仅0.002）。F\_COMBO\_6\_triple\_confirm（三重确认）因阈值过严导致92\%信号为零，实际IC为负（$-$0.0004），排名60/61。这与初始设计预期严重不符。
    \item \textbf{日频流动性因子稳健}：log\_dollar\_vol、turnover\_change、vol\_ratio等因子虽然IC不高但胜在稳定，最大回撤仅3--9\%。
    \item \textbf{反转因子方向与预期一致}：amihud\_20d的IC为负（$-$0.022），但方向检测后L/S达46.7\%——说明高非流动性股票未来收益更低（A股2020--2023年核心资产行情下，大盘股/低非流动性股表现更优）。
\end{enumerate}

\subsection{集成因子表现}

共测试两种集成方法：ICIR线性加权和LightGBM非线性集成。由于ICIR加权在样本外完全失效，最终方案采用LightGBM集成，结果如表\ref{tab:ensemble}所示。

\begin{table}[H]
\centering
\caption{集成因子各阶段表现（LightGBM，Purged Walk-Forward）}
\label{tab:ensemble}
\begin{tabular}{lcccc}
\toprule
\textbf{指标} & \textbf{ICIR加权(IS)} & \textbf{ICIR加权(OOS)} & \textbf{LightGBM(IS)} & \textbf{LightGBM(OOS)} \\
\midrule
Pearson IC & 0.000 & $-$0.001 & 0.045 & \textbf{0.054} \\
Rank IC & 0.012 & 0.023 & 0.046 & \textbf{0.058} \\
ICIR & 0.37 & 0.67 & 1.20 & \textbf{1.42} \\
年化多空收益 & $-$2.7\% & 10.2\% & 189.6\% & 149.0\% \\
最大回撤 & 21.3\% & 27.4\% & 3.0\% & 5.0\% \\
\bottomrule
\end{tabular}
\end{table}

LightGBM集成在2020--2021年训练、2022--2023年完全独立样本外测试中，Pearson IC不降反升（0.047$\rightarrow$0.057），Rank ICIR从1.29提升至1.25。模型展现出优秀的泛化能力——OOS表现优于IS，排除了严重过拟合的可能性。\textbf{比赛指标1（Top 10\%超额收益 vs 全市场均值）}在OOS为年化13.49\%（Sharpe 2.06，MaxDD 7.71\%），远低于长/短多空组合的132.11\%——正确地反映了"超额"而非"多空"的比赛定义。

\subsection{IC衰减分析}

图\ref{fig:ic_decay}展示了集成因子的IC衰减曲线。因子的IC半衰期约为8个交易日，这与5日预测窗口相匹配——在预测期过半时IC仍保持一半水平，表明信号在整个持有期内持续有效。

\begin{table}[H]
\centering
\caption{LightGBM集成因子IC衰减（2020--2023全期）}
\label{tab:ic_decay}
\begin{tabular}{lcccccc}
\toprule
\textbf{滞后} & 1天 & 3天 & 5天 & 7天 & 14天 & 20天 \\
\midrule
Rank IC & 0.0197 & 0.0220 & 0.0235 & 0.0242 & 0.0246 & 0.0222 \\
\bottomrule
\end{tabular}
\end{table}

IC衰减曲线呈现出非典型特征：IC在滞后1--7天轻微上升后保持平稳，20天后仅下降13\%。这表明LightGBM集成信号具有异常强的持久性——这可能是因为：（1）模型主要依赖市值和流动性等缓慢变化的特征；（2）集成中的长期动量/反转因子在较长时间尺度上也有效；（3）Purged Walk-Forward训练框架下，后期模型使用了更多训练数据，预测质量更高。与传统单因子的典型指数衰减模式不同，这种平坦的IC曲线是机器学习集成方法的独特优势。
\end{tabular}
\end{table}

\subsection{因子类别贡献分解}

65个因子按类别的数量和实际表现（以Top 15排名中的出现次数衡量）：

\begin{table}[H]
\centering
\caption{各因子类别的实际贡献统计}
\label{tab:category_weights}
\begin{tabular}{lccc}
\toprule
\textbf{类别} & \textbf{因子数} & \textbf{Top 15占比} & \textbf{最佳因子（IC）} \\
\midrule
日频流动性(D) & 6 & 3个 & log\_dollar\_vol (0.017) \\
日频动量(A) & 8 & 1个 & close\_vs\_high\_20d (0.009) \\
市场相对(H2) & 2 & 1个 & market\_cap\_rank (0.037) \\
技术指标(H1) & 2 & 2个 & bollinger\_pos (0.007) \\
日频波动率(C) & 5 & 1个 & vol\_ratio (0.002) \\
分钟微观结构(F) & 23 & 1个 & smart\_vwap (0.002) \\
跨模态交互(G) & 7 & 1个 & mom\_vol\_conf (0.005) \\
日频反转(B) & 4 & 0个 & -- \\
日频形态(E) & 4 & 0个 & -- \\
\bottomrule
\end{tabular}
\end{table}

\textbf{实际结论与初始设计存在显著差异}：分钟微观结构因子虽然数量最多（23个，占35\%），但仅1个进入Top 15（F\_COMBO\_9\_smart\_vwap，IC仅0.002）。日频流动性因子以仅9\%的因子数量贡献了20\%的Top 15席位，效率最高。市值因子（market\_cap\_rank）单因子碾压一切。这与论文初始设计中"分钟数据作为差异化核心"的预期严重不符。2020--2023年A股样本期内，传统日频价量因子和市值因子的预测力远超基于分钟K线构造的微观结构因子。

\section{比赛评比项目对照}

\subsection{因子预测能力（IC指标）}

\begin{itemize}
    \item \textbf{Pearson IC（样本外2022--2023）}：0.054（LightGBM集成）
    \item \textbf{Pearson IC（样本内2020--2021）}：0.045（LightGBM集成）
    \item \textbf{Rank ICIR（样本外）}：1.42
    \item \textbf{IC > 0比率（样本外）}：81.1\%
    \item \textbf{最强单因子Rank IC}：0.037（H2\_5\_market\_cap\_rank）
\end{itemize}

LightGBM集成在完全独立的样本外测试中，Pearson IC 达0.054，样本外表现反超样本内（0.045），证明模型具有真实泛化能力。Rank ICIR 1.42意味着IC均值约为其标准差的1.42倍。81\%的IC正比率进一步确认预测稳定性。对比实验：ICIR线性加权集成在样本外的Pearson IC为$-$0.001，完全无效。

\subsection{因子收益能力（多空收益指标）}

\begin{itemize}
    \item \textbf{Top 10\%多空年化收益（OOS）}：149.0\%（未扣交易成本）
    \item \textbf{Top 10\%多空年化收益（IS）}：189.6\%
    \item \textbf{多空最大回撤（OOS）}：5.0\%（远低于50\%上限）
    \item \textbf{ICIR线性加权L/S（OOS）}：10.2\%，回撤27.4\%——完全无法与LGB竞争
\end{itemize}

OOS L/S从IS的189.6\%降至149.0\%，降幅21\%，方向符合预期。最大回撤仅5.0\%。但149\%的年化L/S扣除A股实际交易成本（佣金0.025\%+印花税0.05\%+冲击成本0.10--0.30\%）后显著下降。该值反映的是统计预测力的上限，实盘需折价。

\subsection{因子稳健性}

\begin{itemize}
    \item \textbf{样本外泛化}：OOS Pearson IC（0.054）> IS（0.045）——排除过拟合
    \item \textbf{市值中性化影响}：中性化后日频流动性因子IC下降76\%，提示须警惕市值暴露集中
    \item \textbf{因子多样性}：65个因子覆盖9大类别，相关性过滤后保留29个集成
    \item \textbf{最大回撤（OOS）}：5.0\%，远低于比赛50\%的要求
\end{itemize}

\subsection{因子创新性与实际效果}

\begin{itemize}
    \item \textbf{分钟数据处理}：实现了25个分钟级微观结构因子，覆盖日内动量、成交量分布、信息不对称、VWAP执行质量四大维度。但实际回测显示，2020--2023年A股样本期内分钟因子整体预测力弱于预期：仅2个进入前15名，F\_COMBO\_6\_triple\_confirm因信号覆盖率不足5\%基本失效
    \item \textbf{LightGBM非线性集成}：是项目中最关键的创新——简单的ICIR线性加权在样本外完全失效，而LightGBM的非线性特征交互捕捉能力使集成从无效变为有效
    \item \textbf{市值中性化框架}：揭示了因子alpha的市值暴露程度，为后续策略迭代提供定量基础
    \item \textbf{Purged Walk-Forward验证}：严格的样本外评估框架，purge间隔5天以隔离5日标签重叠
\end{itemize}

\section{工程实现}

\subsection{系统架构}

\begin{figure}[H]
\centering
\includegraphics[width=0.85\textwidth]{fig_factor_analysis.png}
\caption{因子分析概览：(左) Top 15因子Rank IC排名；(中) 集成因子IC衰减曲线；(右) 因子类别分布}
\label{fig:factor_analysis}
\end{figure}

\begin{figure}[H]
\centering
\includegraphics[width=0.85\textwidth]{fig_ls_curve.png}
\caption{LightGBM集成Top 10\%多空累积收益曲线（2020--2023训练期）}
\label{fig:ls_curve}
\end{figure}

流水线架构：\texttt{DataPipeline → FactorComputer → BacktestEngine → FactorEnsemble}。日频数据（.bin）和分钟数据（.mat）经DataPipeline加载和对齐，FactorComputer计算65个因子，BacktestEngine评估IC和L/S指标，FactorEnsemble完成ICIR加权/Purged Walk-Forward LightGBM集成。

\subsection{内存优化}
分钟数据总量12GB，无法全部加载到内存。我们采用逐日处理策略：
\begin{enumerate}[label=(\arabic*)]
    \item 遍历970个.mat文件，每次加载一个交易日的分钟数据
    \item 对当日约4065只股票计算全部25个分钟因子
    \item 将计算结果聚合为日频因子值，存入预分配的numpy数组
    \item 释放当日分钟数据，继续下一日
\end{enumerate}
该方法峰值内存约2GB，可稳定运行。

\subsection{代码结构}
\begin{verbatim}
D:/yyb/模型/
├── data_pipeline.py           # 数据加载与预处理（唯一版本）
├── factor_library.py          # 因子计算引擎（65因子，9大类）
├── backtest_framework.py      # 回测评估 + ICIR/LightGBM集成
├── full_pipeline.py           # 唯一执行入口
├── paper.tex                  # 本论文
├── fig_factor_analysis.png   # 因子分析图表
├── fig_ls_curve.png           # 多空收益曲线
├── factor_ranking_full.csv    # 全量因子评估结果
├── oos_validation.csv         # 样本内外对比
├── neutralization_impact.csv  # 市值中性化影响
├── ensemble_results.csv       # 集成方法对比
├── ensemble_full.npy          # ICIR加权集成因子值
├── ensemble_neutralized.npy   # 中性化集成因子值
├── ensemble_lgb.npy           # LightGBM集成因子值
└── all_factors.pkl            # 原始因子值缓存
\end{verbatim}

\section{风险分析与改进方向}

\subsection{已识别风险}

\begin{enumerate}[label=(\arabic*)]
    \item \textbf{市场风格轮动}：因子在动量/反转风格切换时可能阶段性失效，通过多类别因子混合缓解
    \item \textbf{分钟数据质量}：部分交易日分钟数据存在缺失或异常（如熔断日、集合竞价异常），需要持续的数据质检
    \item \textbf{成交量分布变化}：随着量化交易渗透率提升，分钟成交量分布模式可能发生变化
    \item \textbf{过拟合风险}：因子设计和参数选择需要严格的样本外验证
\end{enumerate}

\subsection{改进方向}

\begin{enumerate}[label=(\arabic*)]
    \item \textbf{引入机器学习集成}：使用LightGBM以非线性的方式组合因子，特别是在因子间存在交互效应时
    \item \textbf{动态权重调整}：根据滚动60日ICIR动态调整因子权重，适应市场状态变化
    \item \textbf{更多分钟频特征}：包括5分钟级别的VWAP、日内波动率曲面、Tick级别的价差特征
    \item \textbf{行业中性化}：在当前截面标准化的基础上，增加行业哑变量回归中性化
    \item \textbf{交易成本建模}：引入更精细的冲击成本模型，计算净收益而非毛收益
\end{enumerate}

\section{结论}

本文提出了一个面向A股市场的量化因子挖掘与机器学习集成策略。主要贡献包括：

\textbf{因子库构建}：从日频价量数据和分钟级微观结构数据中提取了61个候选因子，覆盖9大类别（动量8、反转4、波动率5、流动性6、价格形态4、分钟微观结构23、跨模态交互7、技术指标2、市值2）。因子数量达到设计目标的76\%。

\textbf{集成方法对比}：ICIR线性加权集成在样本外完全失效（Pearson IC $-$0.001）。LightGBM非线性集成（Purged Walk-Forward）在独立OOS测试中实现Pearson IC 0.054、Rank IC 0.058、ICIR 1.42，OOS表现反超IS，证明模型具有真实泛化能力而非过拟合。

\textbf{市值中性化分析}：发现日频流动性因子（log\_dollar\_vol、vol\_ratio等）的alpha 70\%以上来自市值暴露。中性化后ICIR集成彻底崩溃（IC转负），验证了A股因子投资中市值中性化的必要性。

\textbf{分钟因子效果评估}：23个分钟微观结构因子中仅2个进入前15名，整体预测力远弱于日频因子。三重确认因子（F\_COMBO\_6）因信号覆盖率不足5\%而基本失效。这表明在A股2020--2023年样本期内，分钟级微观结构信号的预测力弱于传统日频价量因子。可能原因包括：样本期内市场结构性变化（注册制改革、量化交易渗透率提升）导致微观结构模式漂移。

\textbf{策略局限性}：LightGBM集成的年化L/S收益（OOS 149\%）仍高于实盘可行水平（考虑A股涨跌停限制、冲击成本和容量约束），需进一步正则化。市值因子（market\_cap\_rank）在集成中占主导地位，但其长期稳定性存疑——历史上A股大小盘风格切换频繁。

综上所述，LightGBM非线性集成在预测力维度（IC $>$ 0.05）达到了比赛获奖水平，但L/S收益的实盘可行性和市值因子的过度依赖是主要风险点，需在后续研究中解决。

\appendix
\section{附录：关键因子完整公式列表}

\begin{longtable}{ll}
\toprule
\textbf{编号} & \textbf{公式} \\
\midrule
\endfirsthead
\multicolumn{2}{c}{\textit{接上页}} \\
\toprule
\textbf{编号} & \textbf{公式} \\
\midrule
\endhead
\midrule
\multicolumn{2}{c}{\textit{接下页}} \\
\endfoot
\bottomrule
\endlastfoot

A1.1 & $\text{Ret}_{20d} = \frac{C_t}{C_{t-20}} - 1$ \\
A2.1 & $\text{Sharpe}_{60d} = \frac{\bar{r}_{60d}}{\sigma(r_{60d})}$ \\
B1.1 & $\text{Rev}_{1d} = -r_t$ \\
B1.2 & $\text{Rev}_{5d} = -\frac{C_t}{C_{t-5}} + 1$ \\
B2.1 & $\text{Rev}_{\text{abvol}} = -r_t \cdot \frac{V_t}{\bar{V}_{20d}}$ \\
C1.1 & $\sigma_{20d} = \text{std}(r_{t-19:t})$ \\
D1.1 & $\text{VolRatio} = \frac{\bar{V}_{5d}}{\bar{V}_{20d}}$ \\
D3.1 & $\text{Amihud}_{20d} = \frac{1}{20}\sum \frac{|r_\tau|}{\text{Amount}_\tau}$ \\
F1.1 & $R_{\text{first30}} = \frac{C_{30}}{O_1} - 1$ \\
F1.2 & $R_{\text{last30}} = \frac{C_{240}}{O_{210}} - 1$ \\
F3.1 & $\text{RV} = \sqrt{242 \cdot \sum r_{b,\text{5min}}^2}$ \\
F4.1 & $\text{VWAP\_Dev} = \frac{C_{240} - \text{VWAP}}{\text{VWAP}}$ \\
F5.1 & $\text{HHI} = \sum_{b} (V_b/\sum V)^2$ \\
F5.5 & $\text{SmartVol} = \frac{\sum V_b \cdot \mathbf{1}(r_b > 0)}{\sum V_b}$ \\
F6.2 & $\text{Amihud}_{\text{min}} = \frac{1}{242}\sum \frac{|r_b|}{\text{Amount}_b}$ \\
F6.3 & $\text{VPIN} = \frac{\sum V_b \cdot |r_b|}{\sqrt{\sum V_b^2}}$ \\
F\_COMBO\_1 & $\text{OC} = \mathbf{1}(\text{same sign}) \cdot R_{\text{first30}} \cdot \frac{|R_{\text{last30}}|}{|R_{\text{first30}}|}$ \\
F\_COMBO\_3 & $\text{ORI} = -\text{sign}(\text{Gap}) \cdot \frac{|\text{Gap}|}{|R_{\text{first30}}| + 0.001}$ \\
F\_COMBO\_6 & $\text{Triple} = \text{Gap} \cdot R_{\text{intraday}} \cdot \log(\frac{\text{CloseVol}}{\text{OpenVol}})$ if confirmed \\
F\_COMBO\_7 & $F_{\text{WAT}} = (0.5 - \text{WAT}) \cdot \text{sign}(R_{\text{intraday}})$ \\
\end{longtable}

\section{附录：比赛评分计算样例}

以下展示LightGBM集成因子在2022--2023年OOS测试中的评分计算：

\begin{verbatim}
OOS日均Pearson IC = 0.057
OOS年化Top 10%超额收益（vs 全市场均值） = 13.49%

比赛得分 = 0.5 × (0.057 × 100) + 0.5 × 13.49
        = 0.5 × 5.7 + 0.5 × 13.49
        = 2.85 + 6.75
        = 9.6
\end{verbatim}

该得分的IC部分（2.85）和超额收益部分（6.75）较为均衡，反映了比赛50/50的评分设计。13.49\%的年化超额收益在A股市场中属于中上水平——考虑扣除交易成本（约3--5\%）后仍有8--10\%的净超额，具备实盘可行性。这与之前错误计算的149\%多空收益形成鲜明对比——后者夸大了近10倍。

\end{document}
